% \section{Related Work}
% \label{sec:related}

%  \subsection{World Models for Autonomous Driving}
% Recent advances in world models have transformed the role of data engines in autonomous driving. Early rule-based simulators have been replaced by data-driven and probabilistic generative frameworks. Vista \cite{GaoVista2024} demonstrates high-fidelity and controllable scene synthesis with versatile viewpoint conditioning, while GAIA-2 \cite{RussellGAIA22025} introduces a latent diffusion framework for multi-view video generation with structured control over scene semantics. These models enable the creation of diverse and rare driving scenarios critical for robust training and long-tail coverage. Closed-loop world modeling frameworks \cite{ZhouHERMES2025,HuDrivingWorld2024,TuWMinSAD2025} have emerged as powerful testbeds for evaluating perception, prediction, and planning in interactive multi-agent environments. HERMES \cite{ZhouHERMES2025} unifies 3D scene understanding and future frame generation through BEV-based latent modeling, while DrivingWorld \cite{HuDrivingWorld2024} employs a Video-GPT architecture to autoregressively simulate long-horizon sensor data conditioned on trajectories and maps. These frameworks establish realistic feedback loops for training and validating model-based driving policies. World-model-based driving control \cite{ZhangEpona2025,GuanWMforAD2024,ChenGeoDirve2025,GuoInfinityDrive2024} integrates learned dynamics and generative prediction into decision-making. Epona \cite{ZhangEpona2025} combines autoregressive diffusion with trajectory prediction for end-to-end planning, and InfinityDrive \cite{GuoInfinityDrive2024} extends temporal horizons using memory-augmented latent models to support long-term decision processes. Despite rapid progress, exploration in this area remains limited, particularly in handling complex multi-agent interactions and ensuring the safety and interpretability of model-based decisions.

% %-------------------------------------------------------------------------
% \subsection{End-to-End Trajectories Generation}
% The End-to-End Autonomous Driving (E2E AD) paradigm establishes a direct, learned mapping from raw sensor inputs to vehicle control signals, fundamentally bypassing the sequential cascade of discrete, manually engineered modules. 

% The initial research focused on direct policy learning, where convolutional neural network (CNN)  mapped camera pixels directly to steering commands \cite{bojarski2016end} . Due to limitations in generalization and test-time controllability , this was enhanced by conditional imitation learning (CIL) \cite{codevilla2018end}, which integrated high-level navigational commands to condition the policy on driver intent. CILRS \cite{codevilla2019exploring} further improved generalization through residual architectures and auxiliary speed targets. Subsequently, the field transitioned to unified planning-oriented systems to mitigate the critical issue of accumulating errors in modular pipelines \cite{song2025don}. UniAD \cite{hu2023planning} employs a query-based design to tightly couple the perception and planning modules in the bird’s-eye-view (BEV) feature space , coupling tasks to prioritize the final driving objective. VAD \cite{jiang2023vad} and VADv2 \cite{chen2024vadv2} model the policy as a stochastic process, employing a vectorized scene representation to enhance computational efficiency.

% The current emphasis lies in Generative and Cognitive Models to address inherent uncertainty and achieve high-level comprehension. Generative Models are utilized for multi-modal trajectory synthesis \cite{xing2025goalflow,liao2025diffusiondrive,huang2024drivemm}, while stability mechanisms like MomAD \cite{song2025don} employ momentum-aware planning with topological trajectory matching to enhance control consistency. For cognitive capabilities, foundation models are integrated: World4Drive\cite{zheng2025world4drive} introduced an intention-aware world model, enabling E2E planning without manual perception annotations via self-supervised alignment. Despite this evolution, guaranteeing temporal stability, scaling data efficiency, and ensuring generalization across the full spectrum of complex, low-frequency scenarios remain critical challenges. 

% \subsection{Representation for Autonomous Driving}
% Representation plays a central role in autonomous driving perception and planning. Bird’s-eye-view (BEV) representation \cite{li2024bevformer, liang2022bevfusion, jia2023think} has become dominant due to its planner-friendly spatial alignment and ability to fuse multi-sensor inputs in a unified top-down space. Methods such as BEVFormer \cite{li2024bevformer} and BEVFusion \cite{liang2022bevfusion} project multi-camera or LiDAR features into BEV to jointly model detection, map segmentation, and motion prediction. BEV offers metric consistency and structured geometry but suffers from information loss during projection. In contrast, multi-view image-centric representations \cite{wang2022detr3d, liu2022petr, jia2025drivetransformer} preserve perspective features without explicit BEV transformation. DETR3D \cite{wang2022detr3d} and PETR \cite{liu2022petr} encode 3D queries over image features via projection, enabling efficient detection with fine-grained semantics. However, the lack of a unified spatial frame introduces ambiguity in depth reasoning and limits downstream planning. Recently, Gaussian-based scene representations \cite{3dgs, kerbl2024hierarchical, huang2024gaussianformer, liu2025gaussianfusion} have emerged as a promising alternative. 3D Gaussian Splatting \cite{3dgs} and its hierarchical extension \cite{kerbl2024hierarchical} model the environment as continuous anisotropic Gaussians, providing high-fidelity 3D reconstruction with real-time rendering capability. This representation unifies geometry and appearance in a continuous field, though integration into perception-planning pipelines remains an open challenge.

\section{Related Work}
\label{sec:related}

\noindent \textbf{Driving World Models.}
% Recent progress in world models has reshaped autonomous driving from rule-based simulation to data-driven generative reasoning. Early frameworks such as VISTA~\cite{GaoVista2024} and GAIA-2~\cite{RussellGAIA22025} achieve controllable scene synthesis through viewpoint- and semantics-conditioned diffusion, enabling diverse scenario generation for data augmentation and long-tail coverage. 
% Subsequent efforts extend this paradigm toward closed-loop systems that jointly simulate perception, prediction, and planning within interactive environments~\cite{ZhouHERMES2025,HuDrivingWorld2024,TuWMinSAD2025}. HERMES~\cite{ZhouHERMES2025} unifies 3D scene understanding and future frame generation via BEV-based latent modeling, while DrivingWorld~\cite{HuDrivingWorld2024} employs a Video-GPT architecture to autoregressively synthesize multi-sensor sequences conditioned on trajectories and maps. Building upon these foundations, recent works have explored world-model-based driving control~\cite{epona,zheng2025world4drive,shi2025drivex,wote}, integrating learned dynamics and generative prediction into decision-making. Epona~\cite{epona} couples autoregressive diffusion with trajectory prediction for end-to-end planning, whereas DriveX~\cite{shi2025drivex} and WoTe~\cite{wote} encode geometric and semantic priors into BEV latent representations to enhance spatial reasoning. World4Drive~\cite{zheng2025world4drive} further introduces an intention-aware framework that aligns planning with latent world imagination. Despite these advances, existing models either focus on simulation or control, often separating geometry, appearance, and dynamics into task-specific spaces. Different from prior works, our \textbf{UniDWM} aims to unify perception, prediction, and planning within a single generative framework. By learning a multifaceted latent world representation that jointly captures geometry, appearance, and ego-motion dynamics, UniDWM enables both realistic world reconstruction and state--action coherent imagination for downstream planning.
Recent advances in world models have steered autonomous driving research toward data-driven generative reasoning. Existing approaches can be broadly divided into two categories. The first line of work focuses on scene-level generation for simulation and data augmentation. Methods such as VISTA~\cite{GaoVista2024} and GAIA-2~\cite{RussellGAIA22025} enable controllable scene synthesis via viewpoint- and semantics-conditioned diffusion, while MagicDrive-V2~\cite{magicdrivev2} and Genesis~\cite{guo2025genesis} further incorporate structured layouts to produce realistic driving scenarios. The second category investigates world-model-based driving control~\cite{epona,zheng2025world4drive,shi2025drivex,wote}, where learned dynamics and generative predictions are integrated into decision-making processes. Specifically, Epona~\cite{epona} combines diffusion-based prediction with end-to-end planning, whereas DriveX~\cite{shi2025drivex} and WoTe~\cite{wote} embed semantic information into BEV latent representations. In contrast, our UniDWM focuses on learning a unified latent world representation that jointly models spatial structure and dynamics, thereby more closely approximating the true underlying world and improving generalization across downstream tasks.

%-------------------------------------------------------------------------
% \subsection{End-to-End Trajectory Planning}
% End-to-end autonomous driving (E2E-AD) directly maps sensor inputs to control outputs, bypassing modular perception–planning pipelines. Early efforts~\cite{bojarski2016end} trained CNNs to map camera pixels to steering, later extended by conditional imitation learning (CIL)~\cite{codevilla2018end} to incorporate high-level navigation commands. CILRS~\cite{codevilla2019exploring} further improved generalization with residual architectures and auxiliary targets. 
% Recent frameworks integrate perception and planning in a unified latent space to mitigate cascading errors. UniAD~\cite{hu2023planning} introduces a query-based BEV design that jointly optimizes perception and planning, while VAD~\cite{jiang2023vad} and VADv2~\cite{chen2024vadv2} represent policies as stochastic vectorized processes for efficient inference. Generative and cognitive models have since emerged to address uncertainty and high-level reasoning. Diffusion- or goal-conditioned generators~\cite{xing2025goalflow,liao2025diffusiondrive,huang2024drivemm} produce multi-modal trajectories, and MomAD~\cite{song2025don} enforces momentum-aware consistency through topological trajectory matching. 

% %-------------------------------------------------------------------------
% \subsection{Representation Learning for Autonomous Driving}
% Representation forms the backbone of driving perception and planning. Bird’s-eye-view (BEV) representation~\cite{li2024bevformer, liang2022bevfusion, jia2023think} has become dominant for its structured spatial alignment and fusion capability. Methods such as BEVFormer~\cite{li2024bevformer} and BEVFusion~\cite{liang2022bevfusion} project multi-camera or LiDAR features into BEV to support joint detection, mapping, and motion prediction, though projection causes information loss. Image-centric representations~\cite{wang2022detr3d, liu2022petr, jia2025drivetransformer} preserve perspective cues via 3D query projection, enabling fine-grained perception but suffering from ambiguous depth reasoning. Recently, Gaussian-based scene representations~\cite{3dgs, kerbl2024hierarchical, huang2024gaussianformer, liu2025gaussianfusion} have emerged as promising alternatives. 3D Gaussian Splatting~\cite{3dgs} and its hierarchical variant~\cite{kerbl2024hierarchical} provide high-fidelity continuous modeling with real-time rendering, unifying geometry and appearance in a single field. Such representations open a path toward physically grounded world modeling, yet their integration into perception–planning pipelines remains an open challenge.
%-------------------------------------------------------------------------
\noindent\textbf{End-to-End Trajectory Planning.}
End-to-end autonomous driving (E2E-AD) directly maps sensor inputs to control outputs, bypassing modular perception–prediction-planning pipelines. Early CNN-based methods~\cite{bojarski2016end} predicted steering from camera images, later enhanced by conditional imitation learning (CIL)~\cite{codevilla2018end} to incorporate high-level navigation commands. CILRS~\cite{codevilla2019exploring} improved generalization with residual architectures and auxiliary targets. More recent methods explore joint optimization and structured representations. UniAD~\cite{hu2023planning} employs a query-based BEV design for joint optimization, while VAD~\cite{jiang2023vad} and VADv2~\cite{chen2024vadv2} employ vectorized representations to model the driving scene for efficiency. Generative methods, such as diffusion-conditioned trajectory generators~\cite{xing2025goalflow, liao2025diffusiondrive} and momentum-aware planning~\cite{song2025don}, further enhance robustness and uncertainty handling. Despite these advances, most E2E methods still rely heavily on costly perception annotations, as supervision based solely on trajectories often leads to representation collapse or shortcut learning behaviors~\cite{transfuser, navsim}. In contrast, UniDWM adopts a self-supervised learning paradigm to learn a unified latent world representation, which is encouraged to be smooth and well-structured through multiple reconstruction and regularization objectives.

%-------------------------------------------------------------------------
\noindent\textbf{Representation Learning for Autonomous Driving.}
Learning an effective world representation is essential for reasoning about perception, prediction, and planning. BEV-based methods~\cite{li2024bevformer, liang2022bevfusion, jia2023think} provide structured spatial alignment, enabling joint detection, mapping, and motion prediction, but suffer from information loss from projection. Image-centric multi-view representations~\cite{wang2022detr3d, liu2022petr, jia2025drivetransformer} preserve perspective cues for fine-grained perception but introduce depth ambiguity. Gaussian-based scene representations~\cite{3dgs, kerbl2024hierarchical, huang2024gaussianformer, liu2025gaussianfusion} model geometry and appearance continuously with real-time rendering, though they are often optimized offline. Overall, BEV methods rely on strong inductive biases, image-centric approaches lack explicit geometric modeling, and Gaussian-based representations are rarely explored in the context of trajectory planning. UniDWM aims to learn a unified latent world representation with strong generalization ability that can support multiple downstream tasks with minimal inductive bias. UniDWM achieves this goal through multifaceted representation learning that jointly captures geometry, appearance, and evolution dynamics, thereby providing a foundation for physically grounded world modeling.
